\section{Resultados sintéticos}

\subsection{Reconstrucción fuera de muestra}

La PINN obtuvo el menor RMSE en los tres niveles de ruido. El desempeño se mantuvo estable al reducir la SNR, mientras que la MLP se deterioró con mayor intensidad y el GLM mantuvo un error alto asociado a la discrepancia entre la HRF canónica y la dinámica generadora.

\begin{figure}[H]
\centering
\safegraphic[width=0.88\linewidth]{figures/synthetic_signal_rmse.png}
\caption{RMSE de reconstrucción del segundo bloque sintético. Los puntos representan medias y las barras la variabilidad entre instancias.}
\label{fig:synthetic_signal_rmse}
\end{figure}

\begin{table}[H]
\centering
\caption{Desempeño medio de la PINN en datos sintéticos.}
\label{tab:synthetic_pinn_summary}
\begin{tabular}{ccccc}
\toprule
SNR & RMSE prueba & $R^2$ prueba & RMSE HRF & $R^2$ HRF \\
\midrule
10 & $0,0577\pm0,0234$ & 0,9922 & 0,0124 & 0,9923 \\
5  & $0,0655\pm0,0263$ & 0,9900 & 0,0139 & 0,9902 \\
2  & $0,0800\pm0,0428$ & 0,9836 & 0,0181 & 0,9826 \\
\bottomrule
\end{tabular}
\end{table}

\subsection{Recuperación de la HRF}

La MLP reconstruyó adecuadamente los bloques largos, pero presentó errores elevados al evaluarse con el evento estandarizado de 1 s. Esto indica que el ajuste de un patrón de bloques no garantiza la identificación de una respuesta hemodinámica generalizable. La PINN, en cambio, recuperó la HRF con errores cercanos a 0,01--0,02 puntos porcentuales BOLD.

\begin{figure}[H]
\centering
\safegraphic[width=0.88\linewidth]{figures/synthetic_hrf_rmse.png}
\caption{RMSE de recuperación de la HRF sintética.}
\label{fig:synthetic_hrf_rmse}
\end{figure}

\subsection{Recuperación de parámetros}

El parámetro $\alpha$ fue el más estable. Los errores medios de $\epsilon$ aumentaron desde 7,1\% en SNR 10 hasta 10,5\% en SNR 2. $\tau$ presentó la mayor variabilidad, con errores medios entre 9,9\% y 14,3\%.

\begin{figure}[H]
\centering
\safegraphic[width=0.88\linewidth]{figures/synthetic_parameter_error.png}
\caption{Error relativo de los parámetros fisiológicos inferidos por la PINN.}
\label{fig:synthetic_parameter_error}
\end{figure}

\subsection{Comparaciones estadísticas}

Las pruebas de Friedman fueron significativas para las métricas de señal y HRF. Las comparaciones de Wilcoxon con corrección de Holm mostraron ventajas significativas de la PINN frente al GLM y la MLP. La reducción media del RMSE de señal respecto del GLM fue 85,8\%, 84,0\% y 81,0\% para SNR 10, 5 y 2, respectivamente. Frente a la MLP, las reducciones fueron 26,0\%, 35,3\% y 60,7\%.
